\chapter{ХЭРЭГЖҮҮЛЭЛТ БА ТУРШИЛТ}
\begin{sloppy}

\section{Технологийн сонголтын үндэслэл}

Энэхүү дипломын ажлын прототипийг Unity 2022.3 LTS, AR Foundation, C\#, Node.js/Express, Prisma ORM, PostgreSQL зэрэг технологийн стек дээр хэрэгжүүлэв. Энэхүү сонголт нь судалгааны зорилготой шууд уялдаатай. Нэгдүгээрт, Unity болон AR Foundation нь мобайл AR орчны camera passthrough, plane detection, anchor, interaction зэргийг нэг орчинд төвлөрүүлж, 3D объект ба интерфэйсийн логикийг нягт уялдуулах боломж олгодог. Хоёрдугаарт, Node.js/Express нь reward claim, collection retrieval, guest session зэрэг REST API урсгалыг хурдан зохион байгуулахад тохиромжтой. Гуравдугаарт, Prisma ба PostgreSQL нь schema-driven хөгжүүлэлт, entity хоорондын хамаарал, unique constraint-ийг найдвартай удирдах боломж бүрдүүлдэг.

Технологийн сонголтыг хийхдээ ``аль технологи хамгийн өргөн хэрэглээтэй вэ'' гэдгээс илүү ``аль технологи энэ судалгааны гол инженерийн асуудлыг хамгийн тод харуулах вэ'' гэсэн шалгуур хэрэглэсэн. Unity төвтэй хэрэгжилт нь байршлын логик ба AR орчныг нэг кодын урсгалд тайлбарлахад илүү тохиромжтой байсан. Иймээс клиент талын хэрэгжилтийг Unity дотор төвлөрүүлж, системийн архитектурыг Unity клиент + backend гэсэн нэг мөр тайлбарт оруулсан.

\section{Unity клиент ба AR модулийн хэрэгжилт}

\subsection{Клиент талын урсгал}

Unity клиент тал нь аппликейшн эхлэх үед guest session үүсгэх, backend-ээс аяллын цэгүүдийг татах, төхөөрөмжийн байршлыг хянах, collection төлөвийг харуулах үндсэн урсгалыг хариуцна. Энэ зорилгоор `BackendBootstrap`, `LocationTracker`, `TourismSpotManager`, `GameSession`, `ARChestSpawner`, `ARChestInteraction` зэрэг логикийн төв класс ашигласан. `BackendBootstrap` нь апп эхлэх мөчид серверийн холболт, guest login, spot data таталтыг удирдана. `LocationTracker` нь GPS шинэчлэлт болон readiness төлөвийг хянаж, `TourismSpotManager` нь ойролцоох цэгүүдийг Haversine зайгаар тооцно.

\subsection{Байршил шалгах логик}

Клиент талд байршлын шалгалт тасралтгүй ажиллах боловч хэрэглэгчийн интерфэйсийг тогтворгүй болгохгүй байхаар зохион байгуулагдсан. Байршлын шинэчлэлт бүр дээр систем spot жагсаалтыг бүхэлд нь дахин шалгаж, радиуст орсон цэг байвал `currentNearbySpot` төлөвийг шинэчилнэ. GPS савлалтын нөлөөг бууруулахын тулд coordinate smoothing болон hysteresis ашигласан. Ингэснээр радиусын хил дээр олон дахин trigger үүсэх эрсдэл багассан.

\subsection{AR scene ба chest placement}

AR scene эхлэхэд `ARChestSpawner` эхлээд plane detection-ийн үр дүнг ажиглаж, тогтвортой candidate plane олдсон эсэхийг шалгана. Хангалттай нөхцөл бүрдээгүй бол fallback chest position ашиглаж, авдрыг хэрэглэгчийн өмнө, харагдахуйц, interaction хийх боломжтой зайд байрлуулна. Энэ шийдэл нь бодит координатад албадан холбохоос илүү UX-д ээлтэй болох нь туршилтаар харагдсан.

`ARChestInteraction` класс нь авдрыг олон дахин нээхээс хамгаалах `isOpening` ба `isOpened` төлөвтэй. Хэрэглэгч авдрыг товшиход коллайдеруудыг идэвхгүй болгож, нээлтийн анимац тоглуулан, дараа нь reward reveal урсгал ажиллуулна. Ийнхүү визуал урсгал ба backend claim урсгал салангид боловч синхрончлолтой ажилладаг.

\subsection{Collection урсгал}

Reward claim амжилттай болсны дараа `GameSession` болон collection state шинэчлэгдэж, collection дэлгэц дээр тухайн шагнал ``авагдсан'' төлөвөөр харагдана. Collection дэлгэц нь зөвхөн UI жагсаалт биш; хэрэглэгчийн ахиц, дахин ашиглах хүсэл, reward loop-ийн үргэлжлэл болдог. Иймээс collection retrieval нь хэрэгжилтийн чухал интерфейс болж өгсөн.

\section{Backend, API ба өгөгдлийн сангийн хэрэгжилт}

\subsection{Guest session ба spot data}

Backend сервер нь Express framework дээр REST API хэлбэрээр хэрэгжсэн. `POST /auth/guest-login` endpoint нь төхөөрөмжийн танигч дээр үндэслэн хэрэглэгчийг шинээр үүсгэх эсвэл өмнөх session-ийг сэргээх логиктой. Энэхүү шийдэл нь бүрэн бүртгэлийн систем хөгжүүлээгүй prototype нөхцөлд хэрэглэгчийн цуглуулга алдагдахгүй байх практик суурь болсон.

`GET /spots` endpoint нь идэвхтэй аяллын цэгүүдийг тэдгээртэй холбогдсон reward мэдээллийн хамт буцаадаг. Клиент тал энэхүү өгөгдлийг газрын зураг, proximity logic, reward тайлбарт зэрэг ашиглах тул endpoint-ийг нэгдсэн мета мэдээлэл өгөхөөр зохион байгуулсан.

\subsection{Reward claim ба idempotent хамгаалалт}

`POST /spots/:id/claim` endpoint нь системийн хамгийн чухал бизнес дүрмийг хэрэгжүүлдэг. Хүсэлт ирэхэд сервер тухайн хэрэглэгч ба байршлын хослол дээр өмнө нь `UserReward` бүртгэл үүссэн эсэхийг шалгана. Хэрэв бүртгэл байгаа бол дахин мөр үүсгэхгүй, харин already claimed гэсэн төлөв буцаана. Хэрэв бүртгэл байхгүй бол шинэ claim үүсгэн, reward мэдээллийг collection урсгалд дамжуулна. Ийм idempotent логик нь клиент талын давхар даралт, сүлжээний давтагдал, AR interaction-ийн санамсаргүй давталтаас хамгаална.

\subsection{Prisma схем ба өгөгдлийн бүрэн бүтэн байдал}

Өгөгдлийн сангийн схемийг Prisma ашиглан тодорхойлсон. `User`, `TourismSpot`, `Reward`, `UserReward` entity-үүдийн хоорондох холбоо нь reward collection-ийн логикийг шууд илэрхийлнэ. `UserReward` хүснэгтийн `@@unique([userId, tourismSpotId])` хэмжигдэхүүн нь өгөгдлийн сангийн түвшинд бизнес дүрмийг баталгаажуулж, системийн өгөгдлийн бүрэн бүтэн байдлыг хамгаална. Энэ бол prototype-ийн түвшинд ч үйлдвэрлэлийн чанарт ойртсон зөв инженерийн шийдэл юм.

\subsection{Collection retrieval}

`GET /users/:id/rewards` endpoint нь хэрэглэгчийн авсан шагналуудыг байршлын мэдээлэлтэй нь уялдуулан буцаадаг. Ингэснээр collection дэлгэц зөвхөн нэрсийн жагсаалт биш, аль reward-ийг хаанаас авсан, хэдэн ширхэг reward нээгдсэн, ямар цэгүүд үлдсэн зэрэг ахицын мэдээллийг харуулах боломжтой болсон.

\section{Fallback, гүйцэтгэл ба практик туршилт}

\subsection{Fallback стратеги}

Хэрэгжүүлэлтийн явцад дараах fallback механизмыг бодитоор ашигласан.

\begin{itemize}
    \item Backend холболт амжилтгүй үед сүүлийн татсан spot data-г локал кешээс сэргээх.
    \item Plane detection тогтворгүй үед fallback chest placement ашиглах.
    \item Reward зураг эсвэл prefab дутуу үед placeholder дүрслэл үзүүлэх.
    \item Байршлын accuracy хангалтгүй үед AR товчийг шууд идэвхжүүлэхгүй, тайлбар төлөв үзүүлэх.
\end{itemize}

Эдгээр механизм нь системийг ``алдаагүй ажиллуулна'' гэхээс илүү ``алдаа гарсан ч хэрэглэгчийн урсгалыг зохистойгоор авч явна'' гэсэн resilience зарчимд суурилж байгаа юм.

\subsection{Туршилтын сценариуд}

\begin{table}[H]
\centering
\caption{Прототипийн гол туршилтын сценариуд}
\begin{tabular}{|p{4cm}|p{4cm}|p{5cm}|}
\hline
\textbf{Сценари} & \textbf{Хүлээгдэж буй үр дүн} & \textbf{Ажиглалт ба дүгнэлт} \\ \hline
Радиус дотор/гадуур шилжилт & AR товч enter босгоор идэвхжиж, exit босгоор унтрах & Hysteresis нэмснээр төлөв байн байн савлах нь багассан \\ \hline
GPS савлалт & Ойролцоох цэгийн төлөв тогтвортой хадгалагдах & Coordinate smoothing хэрэглэснээр худал trigger буурсан \\ \hline
Plane detection амжилтгүй үе & AR туршлага тасрахгүй, fallback chest гарч ирэх & Fallback placement нь хэрэглэгчийн урсгалыг аврахад үр дүнтэй байсан \\ \hline
Давхар reward claim & Нэг л reward бүртгэгдэх & Unique constraint ба idempotent endpoint хослон ажилласан \\ \hline
Сүлжээ тасарсан үе & Cached spot data-аар газрын зураг, proximity logic хязгаарлагдмал ажиллах & Retrieval боломж хязгаарлагдсан ч үндсэн навигацийн урсгал тасраагүй \\ \hline
\end{tabular}
\end{table}

\subsection{Практик хүндрэлүүд}

Талбайн туршилтын явцад гэрэлтүүлэг муу, текстур багатай гадаргуу, хурдан хөдөлгөөн, GPS accuracy дорой орчинд AR тогтвортой байдал буурч байсан. Энэ нь онолын бүлэгт тайлбарласан SLAM ба LBS-ийн хязгаарлалтууд бодит хэрэгжилт дээр шууд илэрч буйг харуулсан. Ийм нөхцөлд макро ба микро логикийг салгасан зохиомж зөв байсан нь батлагдсан. Өөрөөр хэлбэл GPS-ийн асуудлыг AR placement-т өвлүүлэхгүй байх нь практик туршлагаар зөвтгөгдсөн.

\section{Хэрэгжүүлэлтийн үр дүн ба хязгаарлалт}

Энэхүү хэрэгжилтийн үр дүнд хэрэглэгчийн байршлыг тодорхойлох, аяллын цэгүүдийг серверээс ачаалах, proximity logic ажиллуулах, AR scene эхлүүлэх, виртуал авдар байрлуулах, reward claim баталгаажуулах, collection retrieval хийх үндсэн боломж бүхий prototype систем бүрдсэн. Онолын түвшинд санал болгосон ``радиуст суурилсан идэвхжилт + динамик AR placement + unique claim protection'' гэсэн гол инженерийн шийдэл хэрэгжилтийн түвшинд бодит ажиллах боломжтой болох нь нотлогдов.

Гэсэн хэдий ч прототипийн хэд хэдэн хязгаарлалт байна. Нэгдүгээрт, байршлын нарийвчлал болон AR tracking нь төхөөрөмж ба орчноос хүчтэй хамаарч байна. Хоёрдугаарт, reward контентын менежмент одоогоор backend мета өгөгдлийн түвшинд хязгаарлагдсан бөгөөд тусгай admin panel, analytics, олон байршлын өргөтгөсөн агуулга хараахан хэрэгжээгүй. Гуравдугаарт, сүлжээ тасарсан үед collection retrieval ба claim урсгал бүрэн offline байдлаар шийдэгдээгүй. Иймээс энэхүү прототип нь үйлдвэрлэлийн эцсийн бүтээгдэхүүн бус, харин өргөтгөх боломжтой инженерийн үндэс гэж дүгнэнэ.

\end{sloppy}
